\documentclass[11pt,a4paper]{article}
\usepackage[T1]{fontenc}
\usepackage[utf8]{inputenc}
\usepackage[ngerman]{babel}
\usepackage{amsmath,amssymb,mathtools}
\usepackage{geometry}
\usepackage{enumitem}
\usepackage{xcolor}
\geometry{margin=25mm}
\newcommand{\E}{\mathbb{E}}
\newcommand{\Var}{\operatorname{Var}}
\newcommand{\Unif}{\operatorname{Uniform}}
\newcommand{\Risk}{\mathcal{R}}
\title{Problem Sheet 1 -- Lösung}
\date{}
\begin{document}
\maketitle

\section*{Part 1}
\subsection*{1a) Regressionsproblem}
Features X und Label Y mit jeweiligen Raum:
\begin{align*}
 X&=(U_1,U_2,\ldots,U_N), & X&\in(0,\infty)^N,\\
 Y&=L, &Y&\in(0,\infty),\\
 Z&=X\times Y, & Z&=(X,Y)\sim P_L.
\end{align*}
Dabei besteht die Ausgabe $X$ aus $N$ unabhängig und identisch verteilten
Zufallsvariablen mit $U_i\sim\Unif(0,L)$.

Für den quadratischen Verlust gilt
\[
 \ell(h(x),y)=(h(x)-y)^2.
\]
Das Populationsrisiko lautet
\begin{align*}
 \Risk_L(h)
 &=\E\bigl[(h(X)-Y)^2\bigr]\\
 &=\E\bigl[(h(U_1,\ldots,U_N)-L)^2\bigr].
\end{align*}

\subsection*{1b) Populationsrisiko der Hypothesen}
\paragraph{(1) $h_1$}
Es sei
\[
 h_1(U)=\frac{2}{N}\sum_{i=1}^{N}U_i,
 \qquad \E(U_i)=\frac L2,
 \qquad \Var(U_i)=\frac{L^2}{12}.
\]
Dann gilt
\begin{align*}
 \E[h_1(U)]
 &=\frac{2}{N}\sum_{i=1}^{N}\E(U_i)
 =\frac{2}{N}\cdot N\cdot\frac L2=L.
\end{align*}
Da $h_1$ erwartungstreu (= $L$) ist, entspricht das Risiko der Varianz:
\begin{align*}
 \Risk_L(h_1)
 &=\Var(h_1(U))\\
 &=\frac{4}{N^2}\Var\Bigl(\sum_{i=1}^{N}U_i\Bigr)\\
 &=\frac{4}{N^2}\sum_{i=1}^{N}\Var(U_i)\\
 &=\frac{4}{N^2}\cdot N\cdot\frac{L^2}{12}
 =\boxed{\frac{L^2}{3N}}.
\end{align*}
Für $N=10$ ergibt sich
\[
 \boxed{\Risk_L(h_1)=\frac{L^2}{30}}.
\]

\paragraph{(2) $h_2$}
Setze
\[
 \qquad h_2(U)=M
 \qquad M=\max_{1\le i\le N}U_i,
\]
Damit folgt
\begin{align*}
 \Risk_L(h_2)
 &=\E[(M-L)^2]\\
 &=\E(M^2)-2L\E(M)+L^2.
\end{align*}
Das Maximum ist höchstens $t$, wenn jede einzelne Beobachtung höchstens $t$ ist. Also für eine Beobachtung
\[
 \Pr(U_i\le t)=\frac tL
 \quad (0\le t\le L).
\]
Für alle Beobachtungen
\[
 \qquad
 F_M(t)=\Pr(M\le t)=\left(\frac tL\right)^N
 \quad (0\le t\le L).
\]
Ableiten der Verteilungsfunktion $F_M$ ergibt die Dichte
\[
 f_M(t)=F_M'(t)=\frac{Nt^{N-1}}{L^N},
 \qquad (0\le t\le L).
\]
Zur Risikoberechnung werden $\E(M)$ und $\E(M^2)$ (siehe ausmultiplizierte Loss-Funktion) benötigt:
\begin{align*}
 \E(M)
 &=\int_0^L t f_M(t)\,dt
 =\frac{N}{L^N}\int_0^L t^N\,dt
 =\frac{N}{L^N}\cdot\frac{L^{N+1}}{N+1}
 =\frac{N}{N+1}L,\\[3pt]
 \E(M^2)
 &=\int_0^L t^2 f_M(t)\,dt
 =\frac{N}{L^N}\int_0^L t^{N+1}\,dt
 =\frac{N}{L^N}\cdot\frac{L^{N+2}}{N+2}
 =\frac{N}{N+2}L^2.
\end{align*}
Daraus folgt
\begin{align*}
 \Risk_L(h_2)
 &=\frac{N}{N+2}L^2-2L\frac{N}{N+1}L+L^2\\
 &=L^2\left(\frac{N}{N+2}-\frac{2N}{N+1}+1\right)
 =\boxed{\frac{2L^2}{(N+1)(N+2)}}.
\end{align*}
Für $N=10$ gilt
\[
 \boxed{\Risk_L(h_2)=\frac{2L^2}{11\cdot12}=\frac{L^2}{66}}.
\]

\paragraph{(3) $h_3$}
Setze
\[
 \qquad h_3(U)=\frac{N+1}{N}M
 \qquad M=\max_{1\le i\le N}U_i,
\]
Wie bereits berechnet, gelten
\[
 \E(M)=\frac{N}{N+1}L,
 \qquad \E(M^2)=\frac{N}{N+2}L^2.
\]
Damit
\[
 \E[h_3(U)]
 =\frac{N+1}{N}\E(M)
 =\frac{N+1}{N}\cdot\frac{N}{N+1}L=L.
\]
$h_3$ ist erwartungstreu, weshalb
\[
 \Risk_L(h_3)=\Var(h_3(U)).
\]
Zunächst ergibt sich
\begin{align*}
 \Var(M)
 &=\E(M^2)-[\E(M)]^2\\
 &=\frac{N}{N+2}L^2-\frac{N^2}{(N+1)^2}L^2\\
 &=\frac{NL^2}{(N+2)(N+1)^2}.
\end{align*}
Mit $\Var(aM)=a^2\Var(M)$ folgt
\begin{align*}
 \Risk_L(h_3)
 &=\left(\frac{N+1}{N}\right)^2\Var(M)\\
 &=\left(\frac{N+1}{N}\right)^2
 \frac{NL^2}{(N+2)(N+1)^2}
 =\boxed{\frac{L^2}{N(N+2)}}.
\end{align*}
Für $N=10$ gilt
\[
 \boxed{\Risk_L(h_3)=\frac{L^2}{10\cdot12}=\frac{L^2}{120}}.
\]
Somit gilt für $N=10$:
\[
 \boxed{
 \Risk_L(h_3)=\frac{L^2}{120}
 \le\Risk_L(h_2)=\frac{L^2}{66}
 \le\Risk_L(h_1)=\frac{L^2}{30}}
\]

\section*{Part 2}
\subsection*{2a) Problemstellung}
\[
 X=(U_1,U_2,\ldots,U_{10}),\qquad
 X\in[1,5]^{10},\qquad
 Y=L,\qquad Y\in[1,5].
\]
Hier ist $L\sim\Unif[1,5]$, und bedingt auf $L$ gilt
$U_i\sim\Unif[0,L]$. Somit
\[
 Z=X\times Y,\qquad (X,Y)\sim P,
 \qquad \ell(h(x),y)=(h(x)-y)^2.
\]

\subsection*{2b) Erwartetes Risiko}
Für $N=10$ wurden in Teil 1 die bedingten Risiken berechnet:
\[
 \Risk_L(h_1)=\frac{L^2}{30},\qquad
 \Risk_L(h_2)=\frac{L^2}{66},\qquad
 \Risk_L(h_3)=\frac{L^2}{120}.
\]
Nun wird auch über die Zufallsvariable $L$ gemittelt:
\begin{align*}
 \Risk(h)
 &=\E\bigl[\E[(h(U)-L)^2\mid L]\bigr]\\
 &=\E_L[\Risk_L(h)].
\end{align*}
Da $L\sim\Unif[1,5]$, gilt
\begin{align*}
 \E(L^2)
 &=\frac14\int_1^5 l^2\,dl
 =\frac14\cdot\frac{125-1}{3}
 =\frac{31}{3}.
\end{align*}
Damit ergeben sich
\begin{align*}
 \Risk(h_1)&=\frac1{30}\cdot\frac{31}{3}=\frac{31}{90},\\
 \Risk(h_2)&=\frac1{66}\cdot\frac{31}{3}=\frac{31}{198},\\
 \Risk(h_3)&=\frac1{120}\cdot\frac{31}{3}=\frac{31}{360}.
\end{align*}
Also
\[
 \boxed{\Risk(h_3)<\Risk(h_2)<\Risk(h_1)}.
\]

\subsection*{2c) Stichprobe und empirisches Risiko}
Es sei
\[
 L_i\sim\Unif[1,5],\qquad
 U_{i,1},\ldots,U_{i,10}\mid L_i
 \overset{\mathrm{iid}}{\sim}\Unif[0,L_i].
\]
Weiterhin
\[
 X_i=(U_{i,1},\ldots,U_{i,10}),\qquad Y_i=L_i.
\]
Dabei bezeichnet $m$ die Anzahl unabhängig erzeugter gelabelter Beispiele;
$N=10$ ist die Anzahl der Beobachtungen innerhalb eines Eingabevektors.
Die Stichprobe lautet
\[
 D_m=\bigl((X_1,Y_1),\ldots,(X_m,Y_m)\bigr)\sim P^{\otimes m}.
\]
Für eine Hypothese $h\colon X\to Y$ ist das empirische Risiko
\begin{align*}
 \widehat{\Risk}_m(h)
 &=\frac1m\sum_{i=1}^m\bigl(h(X_i)-Y_i\bigr)^2\\
 &=\frac1m\sum_{i=1}^m
 \bigl(h(U_{i,1},\ldots,U_{i,10})-L_i\bigr)^2.
\end{align*}

\end{document}
